\chapterstar{Glosario de Términos}

{\renewcommand{\arraystretch}{1.45}
\begin{longtable}{@{\hspace{4pt}}p{3.8cm}p{10.2cm}@{\hspace{4pt}}}

\toprule
\textbf{Término} & \textbf{Definición} \\
\midrule
\endfirsthead

\toprule
\textbf{Término} & \textbf{Definición} \\
\midrule
\endhead

\midrule
\multicolumn{2}{r@{\hspace{4pt}}}{\small\textit{Continúa en la siguiente página}} \\
\endfoot

\bottomrule
\endlastfoot

Alerta automatizada &
Notificación generada de forma automática por el sistema cuando una variable supervisada supera un umbral predefinido, como una temperatura fuera del rango permitido o un nivel de inventario inferior al mínimo establecido. \\
\specialrule{0.3pt}{0pt}{0pt}

Arquitectura en capas &
Modelo de diseño de software en el que el sistema se divide en niveles funcionales independientes —como la interfaz de usuario, la lógica de negocio y el almacenamiento de datos— cada uno con responsabilidades específicas y comunicación definida entre ellos. \\
\specialrule{0.3pt}{0pt}{0pt}

Autenticación &
Proceso mediante el cual el sistema verifica la identidad de un usuario antes de concederle acceso. En este proyecto se implementa mediante JSON Web Tokens (JWT), que aseguran que solo usuarios registrados puedan interactuar con la API. \\
\specialrule{0.3pt}{0pt}{0pt}

Borrado lógico &
Técnica de eliminación de registros en la que el elemento no se suprime físicamente de la base de datos, sino que se marca con una fecha de baja. Permite conservar el historial de operaciones vinculadas al registro y reactivarlo en el futuro sin pérdida de información. En este proyecto se aplica a usuarios, roles, sucursales y productos. \\
\specialrule{0.3pt}{0pt}{0pt}

Cadena de frío &
Sistema de conservación que mantiene un producto dentro de un rango de temperatura controlado —típicamente entre 2\,°C y 8\,°C— desde su fabricación hasta su entrega. Algunos medicamentos la requieren de forma obligatoria para conservar su efectividad. \\
\specialrule{0.3pt}{0pt}{0pt}

Despliegue &
Proceso de publicar una aplicación en un servidor para que quede disponible a sus usuarios finales a través de internet. En este proyecto, la interfaz web se despliega en Vercel y el servidor junto con la base de datos en Railway. \\
\specialrule{0.3pt}{0pt}{0pt}

Endpoint &
Dirección específica dentro de una API a la que se envían solicitudes para realizar una operación determinada, como consultar el inventario o registrar una venta. Funciona de forma similar a un número de extensión: cada operación disponible en el servidor tiene su propia dirección única. \\
\specialrule{0.3pt}{0pt}{0pt}

Esquema (base de datos) &
Estructura que define las tablas de una base de datos, las columnas de cada una y las relaciones entre ellas. Funciona como el plano sobre el que se organiza toda la información del sistema; en este proyecto evolucionó mediante doce migraciones durante el desarrollo. \\
\specialrule{0.3pt}{0pt}{0pt}

Firmware &
Programa que se graba en la memoria de un microcontrolador y controla directamente su hardware. En este proyecto, el firmware del ESP32-C5 lee el sensor AHT10, actualiza las pantallas OLED y envía las lecturas al servidor cada 30 segundos. \\
\specialrule{0.3pt}{0pt}{0pt}

Framework &
Estructura de software prediseñada que proporciona las herramientas, convenciones y componentes base para desarrollar una aplicación. En lugar de construir todo desde cero, el desarrollador utiliza el framework como punto de partida y añade la lógica específica del proyecto. \\
\specialrule{0.3pt}{0pt}{0pt}

Humedad relativa &
Porcentaje de vapor de agua presente en el aire respecto al máximo que el aire puede contener a una temperatura determinada. Es una variable crítica para el almacenamiento de medicamentos, ya que valores fuera del rango recomendado degradan su composición química. \\
\specialrule{0.3pt}{0pt}{0pt}

Índice (base de datos) &
Estructura auxiliar que acelera las búsquedas sobre una tabla. Funciona como el índice de un libro: en lugar de recorrer todos los registros uno por uno, la base de datos localiza directamente los que coinciden con la consulta. \\
\specialrule{0.3pt}{0pt}{0pt}

Integridad relacional &
Propiedad que garantiza que los datos relacionados entre distintas tablas se mantengan coherentes; por ejemplo, que ninguna venta haga referencia a un producto inexistente. Se implementa mediante restricciones definidas en la propia base de datos. \\
\specialrule{0.3pt}{0pt}{0pt}

Inventario &
Registro organizado de los productos disponibles en un establecimiento, que incluye información como nombre, cantidad, precio y fecha de caducidad. Su gestión eficiente es fundamental para evitar el desabasto y la acumulación de productos caducados. \\
\specialrule{0.3pt}{0pt}{0pt}

Kárdex &
Sistema manual de registro de existencias en el que cada producto cuenta con una tarjeta individual que documenta las entradas, salidas y saldo disponible. Fue el método de control de inventario utilizado por la Farmacia Selene antes de la implementación del sistema digital. \\
\specialrule{0.3pt}{0pt}{0pt}

Latencia &
Tiempo que transcurre entre el momento en que un sistema envía una solicitud y el momento en que recibe una respuesta. En sistemas de monitoreo ambiental, una latencia baja es deseable para que las alertas lleguen al administrador con la menor demora posible. \\
\specialrule{0.3pt}{0pt}{0pt}

Lote &
Conjunto de unidades de un mismo medicamento fabricadas en una misma producción, identificadas con un código y una fecha de caducidad comunes. El inventario del sistema se gestiona por lote para rastrear qué unidades específicas se vendieron y cuáles deben despacharse primero según su vencimiento. \\
\specialrule{0.3pt}{0pt}{0pt}

Microcontrolador &
Circuito integrado que combina en un solo chip un procesador, memoria y periféricos de entrada/salida programables. Se utiliza para controlar dispositivos electrónicos y ejecutar tareas específicas de forma autónoma. En este proyecto se emplea el ESP32-C5 para la gestión de sensores y la comunicación con el servidor. \\
\specialrule{0.3pt}{0pt}{0pt}

Migración de base de datos &
Procedimiento controlado que modifica la estructura de una base de datos —añadiendo columnas, cambiando tipos de datos o creando índices— sin eliminar los datos existentes. Permite evolucionar el esquema a medida que el sistema crece, conservando la información ya registrada. \\
\specialrule{0.3pt}{0pt}{0pt}

Módulo &
Unidad funcional independiente dentro de un sistema de software o hardware que agrupa elementos relacionados entre sí. La división en módulos facilita el mantenimiento, la comprensión y la ampliación del sistema. \\
\specialrule{0.3pt}{0pt}{0pt}

Nodo sensor &
Conjunto formado por el microcontrolador ESP32-C5, el sensor AHT10 y las tres pantallas OLED, alojado en el gabinete impreso en 3D. Toma lecturas de temperatura y humedad del área de almacenamiento y las transmite al servidor por Wi-Fi. \\
\specialrule{0.3pt}{0pt}{0pt}

ORM &
Object-Relational Mapper — Mapeador Objeto-Relacional. Herramienta que permite interactuar con una base de datos relacional utilizando objetos del lenguaje de programación, sin necesidad de escribir instrucciones SQL directamente. En este proyecto se utiliza SQLAlchemy como ORM. \\
\specialrule{0.3pt}{0pt}{0pt}

Portal cautivo &
Página web que un dispositivo genera y muestra automáticamente cuando alguien se conecta a la red Wi-Fi que él mismo emite. El nodo sensor lo utiliza durante la instalación para que el usuario le indique, desde un teléfono o computadora, a qué red de la farmacia debe conectarse. \\
\specialrule{0.3pt}{0pt}{0pt}

Principio activo &
Sustancia química responsable del efecto terapéutico de un medicamento. Por ejemplo, el principio activo del Paracetamol 500\,mg es el paracetamol; distintas presentaciones comerciales pueden compartir el mismo principio activo con diferente marca o presentación. \\
\specialrule{0.3pt}{0pt}{0pt}

Repositorio &
Almacén centralizado donde se guarda y gestiona el código fuente de un proyecto, incluyendo el historial completo de cambios. En este proyecto se utilizan repositorios en GitHub para el servidor, la interfaz web y el firmware. \\
\specialrule{0.3pt}{0pt}{0pt}

Sensor &
Dispositivo electrónico que detecta una magnitud física del entorno —como temperatura, humedad, presión o luz— y la convierte en una señal eléctrica interpretable por un sistema digital. En este proyecto se utiliza el sensor AHT10 para medir temperatura y humedad relativa. \\
\specialrule{0.3pt}{0pt}{0pt}

Sistema embebido &
Sistema de cómputo diseñado para realizar una función específica dentro de un dispositivo más grande, generalmente con recursos limitados de memoria y procesamiento. El ESP32-C5 programado para gestionar el sensor y la conexión Wi-Fi es un ejemplo de sistema embebido. \\
\specialrule{0.3pt}{0pt}{0pt}

Tablero de indicadores &
Pantalla principal de la interfaz web (también llamada \textit{dashboard}) que concentra en una sola vista los indicadores del negocio: ventas del día, nivel de existencias, lotes próximos a vencer y lecturas ambientales en tiempo real. \\
\specialrule{0.3pt}{0pt}{0pt}

Termohigrómetro &
Instrumento que mide simultáneamente la temperatura y la humedad relativa del ambiente. La normativa sanitaria exige que las farmacias cuenten con uno calibrado y registren sus lecturas; en este proyecto esa función la cumple el nodo sensor de forma automática. \\
\specialrule{0.3pt}{0pt}{0pt}

Token de autenticación &
Cadena de texto generada por el servidor tras una autenticación exitosa, que el cliente incluye en solicitudes posteriores para demostrar su identidad sin necesidad de volver a ingresar credenciales. En este proyecto se implementa mediante el estándar JWT. \\
\specialrule{0.3pt}{0pt}{0pt}

Validación de datos &
Proceso mediante el cual el sistema verifica que la información recibida cumple con el formato, tipo y rango esperados antes de procesarla o almacenarla. Evita errores e inconsistencias en la base de datos originados por datos mal formados o fuera de rango. \\

\end{longtable}}
